\chapter[\hspace{0.6em}IP 加权与边缘结构模型]{IP 加权与边缘结构模型}

第二部分围绕一个因果问题展开：“戒烟对体重增加的平均因果效应是多少？”本章说明如何用 IP 加权从观测数据中估计这一效应。IP 加权在第 2 章已经引入，但当时只把它描述为一种非参数方法。现在我们说明如何把模型与 IP 加权结合起来使用；在附加假设下，这将使我们能够处理协变量众多、处理取值非二值的高维问题。

为估计戒烟对体重增加的影响，我们使用 NHEFS 的真实数据。NHEFS 是一个缩写，全称是（准备好听一个长名字了吗？）National Health and Nutrition Examination Survey Data I Epidemiologic Follow-up Study，即全国健康与营养检查调查数据 I 流行病学随访研究。NHEFS 由美国国家卫生统计中心（National Center for Health Statistics）与美国国家老龄研究所（National Institute on Aging）联合发起，并与美国公共卫生局的其他机构合作。关于 NHEFS 的详细介绍以及公开可用的数据集与文档，见 wwwn.cdc.gov/nchs/nhanes/nhefs/。本章及后续各章使用本书网站上提供的 NHEFS 数据子集。我们鼓励读者改进并完善我们的分析。

\section{因果问题}

我们的目标是估计戒烟（治疗）$A$ 对体重增加（结局）$Y$ 的平均因果效应。为此，我们使用 NHEFS 中 1566 名 25 至 74 岁吸烟者的数据；这些人都接受过一次基线访视，并在约 10 年后接受过一次随访访视。\mnote{分析只纳入以下个体：在基线（1971—75 年）与随访（1982 年）两次访视时性别、年龄、种族、体重、身高、教育、饮酒情况与吸烟强度均已知，且在基线时回答了病史问卷。见精要 12.1。}若某人在随访访视之前报告已戒烟，就归为受治疗 $A=1$；否则归为未受治疗 $A=0$。每个人的体重增加 $Y$（以 kg 计）定义为随访时的体重减去基线时的体重。多数人的体重都增加了，但戒烟者平均增加得更多。戒烟者的平均体重增加为 $\widehat{\E}[Y\given A=1]=4.5$ kg，未戒烟者为 $\widehat{\E}[Y\given A=0]=2.0$ kg。因此，差值 $\E[Y\given A=1]-\E[Y\given A=0]$ 估计为 2.5，95\% 置信区间为 1.7 至 3.4。

\begin{table}[htbp]
\centering
\caption{}
\label{tab:12.1}
\small
\begin{tabular}{lcc}
\toprule
基线特征均值 & \multicolumn{2}{c}{$A$}\\
\cmidrule(lr){2-3}
 & 1 & 0\\
\midrule
年龄，岁 & 46.2 & 42.8\\
男性，\% & 54.6 & 46.6\\
白人，\% & 91.1 & 85.4\\
大学，\% & 15.4 & 9.9\\
体重，kg & 72.4 & 70.3\\
每日吸烟支数 & 18.6 & 21.2\\
吸烟年数 & 26.0 & 24.1\\
很少锻炼，\% & 40.7 & 37.9\\
静态生活，\% & 11.2 & 8.9\\
\bottomrule
\end{tabular}
\end{table}

我们把 $\E[Y^{a=1}]$ 定义为：若总体中所有个体都在随访访视之前戒烟，本可观测到的平均体重增加；把 $\E[Y^{a=0}]$ 定义为：若总体中所有个体都未戒烟，本可观测到的平均体重增加。我们把加法尺度上的平均因果效应定义为 $\E[Y^{a=1}]-\E[Y^{a=0}]$，即“人人都受治疗”与“人人都未受治疗”两种情形下本可观测到的平均体重之差。本章及后续两章主要关心的就是这个因果效应。

本节第一段所估计的关联差 $\E[Y\given A=1]-\E[Y\given A=0]$，一般不同于因果差 $\E[Y^{a=1}]-\E[Y^{a=0}]$。如果戒烟者与未戒烟者在影响体重增加的特征上存在差异，关联差一般就没有因果解释。例如，戒烟者平均比未戒烟者大 4 岁（戒烟者年龄超过 50 岁的可能性比未戒烟者高 44\%），而年长者体重增加得比年轻者少，无论他们戒烟与否。我们说年龄是 $A$ 对 $Y$ 之效应的一个（替代）混杂因子，因此分析需要校正年龄。未校正的估计值 2.5 可能低估了真实的因果效应 $\E[Y^{a=1}]-\E[Y^{a=0}]$。\mnote{精要 7.3 定义了替代混杂因子。}

如表~12.1 所示，戒烟者与未戒烟者在性别、种族、教育、基线体重与吸烟强度等其他变量的分布上也有差异。如果这些变量是混杂因子，分析中也必须对它们加以校正。第 18 章讨论挑选混杂因子的策略。这里我们假定，以下 9 个变量（均在基线时测量）足以校正混杂：性别（0：男，1：女）、年龄（岁）、种族（0：白人，1：其他）、教育（5 类）、吸烟的强度与年限（每日吸烟支数与吸烟年数）、日常生活中的体力活动（3 类）、休闲锻炼（3 类）以及体重（kg）。也就是说，$L$ 表示由 9 个已测协变量组成的向量。\mnote{代码：程序 12.1 计算了本节展示的描述性统计量。}

\renewcommand{\FPname}{精要（Fine Point）}
\begin{finepoint}[树立一个坏榜样]
我们的戒烟例子很方便：它不需要深入的领域知识，数据也是公开的。为这种方便付出的代价之一，是潜在的选择偏倚。若某人 (i) 在 1971—75 年的基线调查时是吸烟者，且 (ii) 在 1982 年的调查中已戒烟，我们就把他归为受治疗 $A=1$。条件 (ii) 意味着纳入研究的个体在基线至 1982 年之间没有死亡、也没有以其他方式失访（否则他们无法回答该调查）。也就是说，我们是以一个在治疗开始（戒烟）之后才发生的事件——回答 1982 年的调查——为条件，把个体选入研究的。如果治疗影响被选入研究的概率，就可能出现第 8 章所述的选择偏倚。（由于不同个体戒烟的时间不同，$A$ 实际上是时变治疗，第二部分将始终忽略这一点。时变治疗在第三部分讨论。）

戒烟随机化实验不会有这个问题。每个人在基线时会被分配到戒烟或不戒烟，因此即便该个体未能参加 1982 年的访视，其治疗组也是已知的。12.6 节将说明如何处理删失或结局数据缺失所致的潜在选择偏倚——这在观察性研究与随机化实验中都可能出现——但本精要所述的情形不同：缺失的是治疗本身的数据。这一选择偏倚可以通过敏感性分析处理，Hernán 等 (2008, 附录 3) 就是这样做的。

选择这个例子，使我们能在自己的分析中展示已发表的观察性数据分析（旨在模拟目标试验）中一个普遍存在的问题：治疗分配与随访开始时的合格性不匹配（Hernán et al. 2016）。虽然为了让分析保持简单，我们决定忽略这个问题，但教学上的便利并不能成为在现实生活中回避这一偏倚的好借口。
\end{finepoint}
\renewcommand{\FPname}{精要}

\section{用建模估计 IP 权重}

IP 加权构造出一个伪总体，其中从协变量 $L$ 指向治疗 $A$ 的箭头被去掉。更确切地说，该伪总体具有以下两条性质：$A$ 与 $L$ 统计独立；伪总体中的均值 $\E_{ps}[Y\given A=a]$ 等于实际总体中的标准化均值 $\sum_l \E[Y\given A=a,L=l]\Pr[L=l]$。即使条件可交换性 $Y^{a}\indep A\given L$ 在实际总体中不成立，这两条性质依然为真（见技术注记 2.3）。现在，如果条件可交换性 $Y^{a}\indep A\given L$ 在实际总体中成立，那么这两条性质就蕴含：(i) $Y^{a}$ 的均值在两个总体中相同；(ii) 伪总体中无条件可交换性（即无混杂）成立；(iii) 实际总体中的反事实均值 $\E[Y^{a}]$ 等于伪总体中的 $\E_{ps}[Y\given A=a]$；(iv) 伪总体中关联即因果。如果需要重温 IP 加权，请重读第 2 章。

通俗地说，构造伪总体的办法，是把每个个体按“接受其实际所接受的那个治疗水平的条件概率”的倒数（reciprocal）加权。治疗 $A$ 的个体特定 IP 权重定义为 $W^{A}=1/f(A\given L)$。对我们的二值治疗 $A$ 而言，IP 权重的分母 $f(A\given L)$ 对戒烟者来说是给定已测混杂因子时戒烟的概率 $\Pr[A=1\given L]$，对未戒烟者来说是给定已测混杂因子时未戒烟的概率 $\Pr[A=0\given L]$。我们只需估计 $\Pr[A=1\given L]$，因为 $\Pr[A=0\given L]=1-\Pr[A=1\given L]$。\mnote{给定治疗取值的条件概率 $\Pr[A=1\given L]$ 称为倾向性评分。关于倾向性评分，第 15 章有更多讨论。}

第 2.4 节中我们用非参数方法估计了 $\Pr[A=1\given L]$：只需数出 $L$ 的每一层中有多少人受治疗（$A=1$），再除以该层的人数。这一计算所需的全部信息，取自一棵经过因果解释的结构树图，该树图有 4 条分支（$L$ 的 2 个取值乘以 $A$ 的 2 个取值）。但是，当数据像本例这样高维、混杂因子众多且其中一些有许多水平时，对 $\Pr[A=1\given L]$ 作非参数估计根本不可行。即便我们愿意把除年龄之外的 8 个混杂因子都重新编码为最多 6 类，我们的树仍会有超过 200 万条分支。如果使用吸烟年限与强度以及体重的实际取值范围，分支还会多出几百万条。当 1566 个个体分布在几百万个层中时，我们无法得到有意义的非参数分层估计。我们必须求助于建模。\mnote{维度灾难在第 10 章引入。}

为了得到 $L$ 所定义的几百万个层中每一层的 $\Pr[A=1\given L]$ 的参数估计，我们为戒烟概率拟合一个 logistic 回归模型，把 9 个混杂因子全部作为协变量纳入。对（准）连续协变量年龄、体重、吸烟强度与吸烟年限，我们使用了线性项与二次项，并且没有纳入协变量之间的乘积项。也就是说，我们的模型对 $\Pr[A=1\given L]$ 的可能取值施加了这样的限制：在对数优势（logit）尺度上，连续协变量与戒烟风险之间的条件关系可以用一条抛物线表示，且每个协变量对（logit）风险的贡献与其他协变量的贡献无关。在这些参数限制下，我们能够对 $L$ 取值的每一种组合得到 $\Pr[A=1\given L]$ 的估计，从而对研究总体中 1566 个个体中的每一个都得到一个估计值。\mnote{代码：程序 12.2。估计的 IP 权重 $W^{A}$ 从 1.05 到 16.7，均值为 2.00。}

下一步是计算由估计的 IP 权重所构造的伪总体中的差值 $\widehat{\E}_{ps}[Y\given A=1]-\widehat{\E}_{ps}[Y\given A=0]$。如果伪总体中 $A$ 的效应没有混杂，且 $\Pr[A=1\given L]$ 的模型正确，那么关联即因果，伪总体中关联差 $\E_{ps}[Y\given A=1]-\E_{ps}[Y\given A=0]$ 的一个无偏估计量，同时也是实际总体中因果差 $\E[Y^{a=1}]-\E[Y^{a=0}]$ 的无偏估计量。

我们估计伪总体中 $\E_{ps}[Y\given A=1]-\E_{ps}[Y\given A=0]$ 的做法，是用加权最小二乘拟合（饱和的）线性均值模型 $\E[Y\given A]=\theta_0+\theta_1A$，权重为个体估计的 IP 权重 $\widehat{W}$：戒烟者为 $1/\widehat{\Pr}[A=1\given L]$，未戒烟者为 $1/\bigl(1-\widehat{\Pr}[A=1\given L]\bigr)$。参数估计值 $\hat\theta_1$ 为 3.4。也就是说，我们估计戒烟平均使体重增加 $\hat\theta_1=3.4$ kg。该估计量的正式定义见技术注记 12.1。\mnote{$\E[Y\given A]=\theta_0+\theta_1A$ 是饱和模型，因为它有 2 个参数 $\theta_0$ 与 $\theta_1$，用来估计 2 个量：$\E[Y\given A=1]$ 与 $\E[Y\given A=0]$。在这个模型中，$\theta_1=\E[Y\given A=1]-\E[Y\given A=0]$。}

\renewcommand{\TPname}{技术注记（Technical Point）}
\begin{techpoint}[Horvitz-Thompson 估计量]
技术注记 3.1 中，我们定义了治疗水平 $a$ 的“表观”IP 加权均值 $\E\bigl[I(A=a)Y/f(A\given L)\bigr]$，在正值性与可交换性下它等于反事实均值 $\E[Y^{a}]$。这一 IP 加权均值可由原始的 Horvitz-Thompson (1952) 估计量 $\widehat{\E}\bigl[I(A=a)Y/f(A\given L)\bigr]$ 一致地估计，其中 $\widehat{\E}$ 为样本平均算子，且 $f(A\given L)$ 假定为已知。然而本章我们是通过 IP 加权最小二乘估计 $\hat\theta_0+\hat\theta_1a$ 来估计 $\E[Y^{a}]$ 的；对二值 $A$，这是一个修正的 Horvitz-Thompson 估计量，通常称为 Hajek 估计量
\[
\frac{\widehat{\E}\left[\dfrac{I(A=a)Y}{f(A\given L)}\right]}{\widehat{\E}\left[\dfrac{I(A=a)}{f(A\given L)}\right]}
\qquad\text{(Hajek 1971).}
\]

Hajek 估计量是
\[
\frac{\E\left[\dfrac{I(A=a)Y}{f(A\given L)}\right]}{\E\left[\dfrac{I(A=a)}{f(A\given L)}\right]}
\]
的（渐近）无偏估计量；在正值性下，它等于 $\E[Y^{a}]$，因为 $\E\bigl[I(A=a)/f(A\given L)\bigr]=1$。在实践中，Hajek 估计量更受青睐，因为与 Horvitz-Thompson 估计量不同，对二值 $Y$ 它保证落在 0 与 1 之间，即便 $f(A\given L)$ 未知、被替换为由设定有误的模型拟合得到的预测值 $\hat f(A\given L)$ 时也是如此。

另一方面，若正值性不成立，则该比值等于 $\sum_l \E[Y\given A=a,L=l,L\in Q(a)]\Pr[L=l\given L\in Q(a)]$；若可交换性成立，它等于 $\E[Y^{a}\given L\in Q(a)]$，其中 $Q(a)=\{l;\Pr(A=a\given L=l)>0\}$ 是使 $A=a$ 能以正概率被观测到的 $l$ 取值集合。因此，正如技术注记 3.1 所讨论的，在缺乏正值性时，$a=1$ 与 $a=0$ 的 Hajek 估计量之差没有因果解释。在非正值性下，Horvitz-Thompson 估计量的极限与 Hajek 估计量的极限之比不再是 1，而是 $\Pr[Q(a)]$，因为 Hajek 估计量的分母收敛到 $\Pr[Q(a)]$ 而非 1。
\end{techpoint}
\renewcommand{\TPname}{技术注记}

为得到点估计 $\hat\theta_1=3.4$ 周围的 95\% 置信区间，我们需要一种把 IP 加权考虑在内的方法。一种可能是用统计理论推导相应的方差估计量。这种做法要求数据分析者自己编写估计量的程序，而标准统计软件一般并不提供。第二种可能是用非参数自助法（bootstrapping）近似方差。对大型数据库而言，这种做法需要相当的计算资源，或者需要极大的耐心。第三种可能是使用稳健方差估计量（例如用于工作相关结构为独立的 GEE 模型的那种），它是多数统计软件包的标准选项。基于稳健方差估计量的 95\% 置信区间是有效的，但与上述解析方法和自助法估计量不同，它是保守的——它们覆盖超总体参数的频率超过 95\%。$\hat\theta_1$ 周围的保守 95\% 置信区间为 (2.4, 4.5)。本章中，IP 加权估计的所有置信区间都是保守的。如果 $\Pr[A=1\given L]$ 的模型设定有误，$\theta_0$ 与 $\theta_1$ 的估计就会有偏，而且正如上一章所讨论的，置信区间覆盖真实值的频率可能低于 95\%。\mnote{以权重 $\widehat{W}$ 求 $\theta_0$ 与 $\theta_1$ 的加权最小二乘估计，就是使 $\sum_i \widehat{W}_i\bigl(Y_i-(\theta_0+\theta_1A_i)\bigr)^2$ 最小化的 $\hat\theta_0$ 与 $\hat\theta_1$。若对所有个体 $i$ 都有 $\widehat{W}_i=1$，就得到上一章所述的普通最小二乘估计。}

\section{稳定化 IP 权重}

IP 加权的目标是构造一个伪总体，其中协变量 $L$ 与治疗 $A$ 之间没有关联。第 2 章展示了如何用 IP 权重 $W^{A}=1/f(A\given L)$ 把图~2.1 中的原始研究总体变换为图~2.3 中的伪总体。该伪总体的规模是原始研究总体的两倍，这反映了权重 $W^{A}$ 的平均值为 2 这一事实。通俗地说，这些权重模拟出一个由原始研究总体的两份拷贝构成的伪总体：一份受治疗，另一份未受治疗。

不过，构造 $A$ 与 $L$ 独立的伪总体还有其他办法。例如，可以构造这样一个伪总体：无论 $L$ 取何值，所有个体接受 $A=1$ 的概率都等于 0.5，接受 $A=0$ 的概率也等于 0.5。这样的伪总体用 IP 权重 $0.5/f(A\given L)$ 构造。它的规模与研究总体相同；若把所有 IP 权重都除以 2，它在代数上就等于上一段那个伪总体。因此，权重 $0.5/f(A\given L)$ 的期望均值为 1，用 $0.5/f(A\given L)$ 构造的伪总体中得到的效应估计，等于用 $1/f(A\given L)$ 构造的伪总体中得到的效应估计。（你可以用图~2.1 的数据经验地验证这一点，或见技术注记 12.2 的证明。）对任何 $0<p\le 1$ 的其他 IP 权重 $p/f(A\given L)$，情况也是如此。权重 $W^{A}=1/f(A\given L)$ 只是 $p=1$ 时 IP 权重的一个特例。

让我们把论证再推进一步。校正混杂的关键要求是：在伪总体中，治疗 $A$ 的概率不依赖于混杂因子 $L$。只要给伪总体中每个人分配相同的治疗概率 $p$，就能满足这一要求。但我们也可以构造一个不同的人有不同治疗概率的伪总体，只要治疗概率不依赖于 $L$ 的取值即可。例如，一种常见的选择是：给受治疗者分配原总体中受治疗的概率 $\Pr[A=1]$，给未受治疗者分配原总体中未受治疗的概率 $\Pr[A=0]$。于是 IP 权重对受治疗者为 $\Pr[A=1]/f(A\given L)$，对未受治疗者为 $\Pr[A=0]/f(A\given L)$；更紧凑地写，就是 $f(A)/f(A\given L)$。\mnote{受治疗子总体中的平均因果效应，可以用分子为 $\Pr[A=1\given L]$ 的 IP 权重来估计。见技术注记 4.1。}

图~12.1 展示了把 IP 权重 $f(A)/f(A\given L)$ 施加于图~2.1 的数据时所构造的伪总体，其中 $\Pr[A=1]=13/20=0.65$，$\Pr[A=0]=7/20=0.35$。在第 3 章的可识别条件下，这个伪总体类似于一个假想的随机化实验：研究总体中 65\% 的个体被随机分配到 $A=1$，35\% 被分配到 $A=0$。注意，为保持 65/35 的比例，各分支中的个体数不可能是整数。幸运的是，在数学上处理“非整数个人”并不算什么难事。

\begin{figure}[htbp]
\centering
\begin{tikzpicture}[causal,scale=0.85,every node/.style={transform shape}]
  % 四条分支
  \draw[->,semithick,shorten <=0.72cm,shorten >=1.00cm] (0,0) -- (5.9,3.80);
  \draw[->,semithick,shorten <=0.72cm,shorten >=1.00cm] (0,0) -- (5.9,1.25);
  \draw[->,semithick,shorten <=0.72cm,shorten >=1.00cm] (0,0) -- (5.9,-1.25);
  \draw[->,semithick,shorten <=0.72cm,shorten >=1.00cm] (0,0) -- (5.9,-3.80);
  \node[rotate=32.8,font=\scriptsize] at (2.95,2.16) {未受治疗（0.35）};
  \node[rotate=32.8,font=\scriptsize\bfseries] at (3.21,1.76) {2.8};
  \node[rotate=12.0,font=\scriptsize] at (3.02,0.86) {受治疗（0.65）};
  \node[rotate=12.0,font=\scriptsize\bfseries] at (3.12,0.40) {5.2};
  \node[rotate=-12.0,font=\scriptsize] at (3.02,-0.86) {未受治疗（0.35）};
  \node[rotate=-12.0,font=\scriptsize\bfseries] at (3.12,-0.40) {4.2};
  \node[rotate=-32.8,font=\scriptsize] at (2.95,-2.16) {受治疗（0.65）};
  \node[rotate=-32.8,font=\scriptsize\bfseries] at (3.21,-1.76) {7.8};
  % 根节点
  \draw[thick,fill=white] (0,0) circle (0.88);
  \draw (-0.88,0) -- (0.88,0);
  \node[font=\scriptsize] at (0,0.35) {$L=0$ (0.4)};
  \node[font=\scriptsize] at (0,-0.35) {$L=1$ (0.6)};
  % 结局节点
  \foreach \y/\a/\b/\p/\q in {3.80/.75/.25/2.1/0.7, 1.25/.75/.25/3.9/1.3,
                              -1.25/1/3/1.4/2.8, -3.80/1/3/2.6/5.2}{
    \draw[thick,fill=white] (5.9,\y) circle (1.05);
    \draw (4.85,\y) -- (6.95,\y);
    \node[font=\scriptsize] at (5.9,\y+0.46) {$Y=0$ (\a)};
    \node[font=\scriptsize] at (5.9,\y-0.46) {$Y=1$ (\b)};
    \node[font=\scriptsize,anchor=west] at (7.05,\y+0.44) {\p};
    \node[font=\scriptsize,anchor=west] at (7.05,\y-0.44) {\q};
  }
  % 稳定化权重列
  \node[font=\scriptsize,anchor=west] at (7.85,4.24) {$0.35/0.5=0.7$};
  \node[font=\scriptsize,anchor=west] at (7.85,3.36) {$0.35/0.5=0.7$};
  \node[font=\scriptsize,anchor=west] at (7.85,1.69) {$0.65/0.5=1.3$};
  \node[font=\scriptsize,anchor=west] at (7.85,0.81) {$0.65/0.5=1.3$};
  \node[font=\scriptsize,anchor=west] at (7.85,-0.81) {$0.35/0.25=1.4$};
  \node[font=\scriptsize,anchor=west] at (7.85,-1.69) {$0.35/0.25=1.4$};
  \node[font=\scriptsize,anchor=west] at (7.85,-3.36) {$0.65/0.75=0.87$};
  \node[font=\scriptsize,anchor=west] at (7.85,-4.24) {$0.65/0.75=0.87$};
  \node[font=\scriptsize,anchor=west] at (7.85,5.10) {$SW^{A}=f(A)/f(A\given L)$};
  \draw (7.85,4.84) -- (10.45,4.84);
\end{tikzpicture}
\caption{}
\label{fig:12.1}
\end{figure}

在我们戒烟的例子中，IP 权重 $f(A)/f(A\given L)$ 从 0.33 到 4.30，而 IP 权重 $1/f(A\given L)$ 从 1.05 到 16.7。分子中的稳定化因子 $f(A)$ 正是 $f(A)/f(A\given L)$ 权重取值范围更窄的原因。IP 权重 $W^{A}=1/f(A\given L)$ 称为\emph{非稳定化权重}，IP 权重 $SW^{A}=f(A)/f(A\given L)$ 称为\emph{稳定化权重}。稳定化权重的均值预期为 1，因为伪总体的规模等于研究总体的规模。

现在用稳定化 IP 权重 $SW^{A}$ 重新估计戒烟对体重的影响。\mnote{在数据分析中，应检查估计的权重 $SW^{A}$ 的均值是否为 1（Hernán and Robins 2006a）。偏离 1 说明模型设定有误，或正值性被违背、接近被违背。关于检查正值性，见精要 12.2。}首先，我们需要 $\Pr[A=1\given L]$ 的估计，用来构造权重的分母。我们使用第 12.2 节用过的同一个 logistic 模型，对研究总体中 1566 个个体逐一得到 $\Pr[A=1\given L]$ 的参数估计 $\widehat{\Pr}[A=1\given L]$。其次，我们需要估计权重分子中的 $\Pr[A=1]$。可以用比值 403/1566 得到非参数估计，等价地，也可以拟合一个只含截距、不含协变量的饱和 logistic 模型来估计 $\Pr[A=1]$。最后，我们用个体估计的稳定化 IP 权重拟合均值模型 $\E[Y\given A]=\theta_0+\theta_1A$，以估计因果差 $\E[Y^{a=1}]-\E[Y^{a=0}]$：戒烟者的权重为 $\widehat{\Pr}[A=1]/\widehat{\Pr}[A=1\given L]$，未戒烟者的权重为 $\bigl(1-\widehat{\Pr}[A=1]\bigr)/\bigl(1-\widehat{\Pr}[A=1\given L]\bigr)$。在我们的假设下，我们估计戒烟平均使体重增加 $\hat\theta_1=3.4$ kg（95\% 置信区间：2.4, 4.5）。这与先前用非稳定化 IP 权重 $W^{A}$ 而非稳定化 IP 权重 $SW^{A}$ 得到的估计相同。\mnote{代码：程序 12.3。估计的 IP 权重 $SW^{A}$ 从 0.33 到 4.30，均值为 1.00。}

如果非稳定化权重与稳定化权重给出相同的估计，那为什么还要用稳定化权重？因为稳定化权重得到的 95\% 置信区间通常比非稳定化权重更窄。不过，稳定化权重在统计上的优越性，只有在（IP 加权）模型不饱和时才会出现。在上面这个例子中，两参数模型 $\E[Y\given A]=\theta_0+\theta_1A$ 是饱和的，因为治疗 $A$ 只能取 2 个可能值。在许多情形中（例如时变治疗或连续治疗），加权模型不可能是饱和的，因此要使用稳定化权重。下一节介绍稳定化权重在连续治疗中的用法。

\renewcommand{\FPname}{精要（Fine Point）}
\begin{finepoint}[检查正值性]
在我们的研究中，有 4 位 66 岁的白人女性，其中无人戒烟。也就是说，以 $L$ 的（某个子集）为条件时 $A=1$ 的概率为 0。IP 加权的一个条件——正值性——在经验上被违背了。正值性可能以两种方式被违背：
\begin{itemize}
\item \textbf{结构性违背}：第 3 章所述的那类违背。具有某些 $L$ 取值的个体不可能受治疗（或不可能未受治疗）。一例：在估计接触某些化学品对死亡率的效应时，离岗是一个重要的混杂因子，因为离岗者更可能生病、死亡，而离岗又是化学品暴露的一个决定因素——人们只有在工作时才可能接触该化学品。也就是说，问题的结构保证了在“离岗”条件下受治疗的概率恰为 0（一个结构性零）。以该混杂因子为条件时，我们总会发现零格。
\item \textbf{随机性违背}：本精要第一段所述的那类违背。我们的样本是有限的，因此若以若干混杂因子分层，即便目标总体中受治疗的概率并非真的为零，我们也会在某些地方开始发现零格。这是对正值性的随机（而非结构性）违背，因为零格在不同样本中随机出现在不同位置。一例：我们的研究恰好在上层“66 岁白人女性”与“67 岁白人女性”中没有受治疗的个体，但在“65 岁白人女性”与“69 岁白人女性”这两层中受治疗的个体数为正。
\end{itemize}
每类正值性违背各有不同的后果。存在结构性违背时，无法用 IP 加权或标准化对整个总体作因果推断。推断必须限制在结构性正值性成立的层内。细节见技术注记 12.1。存在随机性违背时，我们用参数模型、借助其他层中个体的数据来估计随机零格所在层的治疗概率。换句话说，我们用参数模型把零格平滑掉。例如，12.2 节所用的 logistic 模型通过从研究中所有其他个体内插，估计了 66 岁白人女性的戒烟概率。每当我们像估计 $\hat\theta_1=3.4$ 时那样，在存在零格的情况下用参数方法估计 IP 权重，我们实际上都在假定随机性非正值性。
\end{finepoint}
\renewcommand{\FPname}{精要}

\section{边缘结构模型}

考虑以下关于治疗水平 $a$ 之下结局均值的线性模型
\[
\E[Y^{a}]=\beta_0+\beta_1a
\]
这个模型与我们迄今描述的所有模型都不同：它的结局变量是反事实的，因而一般是观测不到的。因此，这个模型无法拟合任何真实世界研究的数据。关于反事实结局边缘均值的模型，称为\emph{边缘结构均值模型}。\mnote{这是针对二值治疗 $A$ 的（饱和）边缘结构均值模型。}

结构均值模型中针对治疗的参数，对应平均因果效应。在上面的模型中，参数 $\beta_1$ 等于 $\E[Y^{a=1}]-\E[Y^{a=0}]$，因为在 $a=0$ 时 $\E[Y^{a}]=\beta_0$，在 $a=1$ 时 $\E[Y^{a}]=\beta_0+\beta_1$。前面几节中，我们估计了戒烟 $A$ 对体重变化 $Y$ 的平均因果效应，其定义为 $\E[Y^{a=1}]-\E[Y^{a=0}]$。换句话说，我们估计的是某个边缘结构模型的参数 $\beta_1$。

具体地说，我们用 IP 加权构造了一个伪总体，然后用 IP 加权最小二乘把模型 $\E[Y\given A]=\theta_0+\theta_1A$ 拟合到伪总体数据上。在我们的假设下，伪总体中关联即因果。也就是说，IP 加权关联模型 $\E[Y\given A]=\theta_0+\theta_1A$ 中的参数 $\theta_1$，可以被赋予与结构模型 $\E[Y^{a}]=\beta_0+\beta_1a$ 中参数 $\beta_1$ 相同的因果解释。由此可知，伪总体中关联参数的一致估计量 $\hat\theta_1$，也是总体中因果效应 $\beta_1=\E[Y^{a=1}]-\E[Y^{a=0}]$ 的一致估计量。

边缘结构模型 $\E[Y^{a}]=\beta_0+\beta_1a$ 是饱和的，因为戒烟 $A$ 是二值治疗。也就是说，方程两边各有 2 个未知量：左边是 $\E[Y^{a=1}]$ 与 $\E[Y^{a=0}]$，右边是 $\beta_0$ 与 $\beta_1$。因此，在伪总体中计算样本均值，就足以估计我们所关心的因果效应。\mnote{边缘结构模型的理想性质之一是空值保持（见第 9 章）：无平均因果效应的零假设为真时，它永不设定有误。此时 $\E[Y^{a}]=\beta_0+\beta_1a+\beta_2a^2$ 中 $\beta_1=\beta_2=0$、$\beta_0=\E[Y^{a}]$；若可交换性成立，则 $\E[Y]=\beta_0$。}但治疗常常是多值的或连续的。例如，考虑一个新治疗 $A$“吸烟强度的变化”，其定义为 1982 年每日吸烟支数减去基线时每日吸烟支数。治疗 $A$ 现在可以取许多值，比如某个体每日吸烟支数减少了 25，则 $A=-25$；增加了 40，则 $A=40$。假设我们想估计在基线时每日吸烟不超过 25 支的 1162 个个体中，治疗强度不同变化之下的平均体重变化之差。也就是说，我们想对任意取值 $a$ 与 $a'$ 估计 $\E[Y^{a}]-\E[Y^{a'}]$。

由于治疗 $A$ 可以取几十个值，参数同样多的饱和模型就不实用了。我们必须考虑一个非饱和的结构模型，以刻画治疗 $A$ 对平均结局 $Y$ 之效应的剂量-反应曲线。如果我们相信抛物线能恰当地描述这条剂量-反应曲线，就会提出边缘结构模型
\[
\E[Y^{a}]=\beta_0+\beta_1a+\beta_2a^2
\]
其中 $a^2=a\times a$ 为 $a$ 的平方，$\E[Y^{a=0}]=\beta_0$ 是 $a=0$ 时的平均体重增加，即在基线至 1982 年间吸烟强度不变时的平均体重增加。\mnote{针对连续治疗 $A$ 的（非饱和）边缘结构均值模型。}

假设我们想估计把吸烟强度增加 20 支/日相对于不变的平均因果效应，即 $\E[Y^{a=20}]-\E[Y^{a=0}]$。\marginnote{更稳定地估计 IP 权重的方法是一个活跃的研究领域。见 Imai and Ratkovic (2015)、Wang and Zubizarreta (2020)、Kallus and Santacatterina (2018) 以及 Avagyan and Vansteelandt (2021) 的工作。}[0.55cm]根据我们的结构模型，$\E[Y^{a=20}]=\beta_0+20\beta_1+400\beta_2$，因此 $\E[Y^{a=20}]-\E[Y^{a=0}]=20\beta_1+400\beta_2$。现在我们需要估计参数 $\beta_1$ 与 $\beta_2$。为此，需要估计 IP 权重 $SW^{A}$ 以构造一个不受 $L$ 混杂的伪总体，然后把关联模型 $\E[Y\given A]=\theta_0+\theta_1A+\theta_2A^2$ 拟合到伪总体数据上。

为估计稳定化权重 $SW^{A}=f(A)/f(A\given L)$，我们需要估计 $f(A\given L)$。对二值治疗 $A$，$f(A\given L)$ 是一个概率，所以我们用 logistic 模型估计 $\Pr[A=1\given L]$。对连续治疗 $A$，$f(A\given L)$ 是概率密度函数（pdf）。遗憾的是，pdf 一般很难估计，尤其当 $L$ 高维且含有连续成分时；这正是用 IP 加权处理连续治疗往往很危险的原因。在我们的例子中，我们假定密度 $f(A\given L)$ 为正态（高斯）分布，均值为 $\mu_L=\E[A\given L]$、方差 $\sigma^2$ 为常数。然后我们用线性回归模型对 $L$ 取值的所有组合估计均值 $\E[A\given L]$ 与残差方差 $\sigma^2$。我们还假定分子中的密度 $f(A)$ 也是正态的。用 IP 加权处理连续治疗时应格外小心，因为效应估计可能对用于估计条件密度 $f(A\given L)$ 的模型或算法的选择极为敏感。\mnote{代码：程序 12.4。估计的 $SW^{A}$ 从 0.19 到 5.10，均值为 1.00。我们假定方差为常数（同方差），检查残差图后看是合理的。}

我们对边缘结构模型参数的 IP 加权估计为 $\hat\beta_0=2.005$、$\hat\beta_1=-0.109$、$\hat\beta_2=0.003$。根据这些估计，若所有个体的吸烟强度保持不变，平均体重增加（95\% 置信区间）本为 2.0 kg（1.4, 2.6）；若所有个体在基线至 1982 年间把吸烟量增加 20 支/日，则本为 0.9 kg（$-1.7$, 3.5）。因此，$\E[Y^{a=20}]-\E[Y^{a=0}]$ 的估计为 $0.9-2.0=-1.1$ kg。

也可以考虑针对二值结局的边缘结构模型。例如，若关心戒烟 $A$（1：是，0：否）对 1992 年前死亡风险 $D$（1：是，0：否）的因果效应，可以考虑如下的边缘结构 logistic 模型
\[
\mathrm{logit}\ \Pr[D^{a}=1]=\alpha_0+\alpha_1a
\]
其中 $\exp(\alpha_1)$ 是戒烟相对于不戒烟的死亡因果优势比。在我们的假设下，把 logistic 模型 $\mathrm{logit}\ \Pr[D=1\given A]=\theta_0+\theta_1A$ 拟合到 IP 加权构造的伪总体上，即可一致地估计该模型的参数。我们估计的因果优势比为 $\exp\bigl(\hat\theta_1\bigr)=1.0$（95\% 置信区间：0.8, 1.4）。\mnote{这是针对二值治疗的饱和边缘结构 logistic 模型。对连续治疗，我们会指定一个非饱和的 logistic 模型。代码：程序 12.5。}

\section{效应修饰与边缘结构模型}

当目标参数是总体中的平均因果效应时，边缘结构模型不包含协变量。然而，人们可以在边缘结构模型中加入协变量——这些协变量可以不是混杂因子——以评估效应修饰。假设有人提出假说：戒烟的效应因性别 $V$（0：男，1：女）而异。为检验这一假说，我们把协变量 $V$ 加入边缘结构均值模型：
\[
\E[Y^{a}\given V]=\beta_0+\beta_1a+\beta_2Va+\beta_3V
\]
若 $\beta_2\neq 0$，则存在加法效应修饰。严格说来，这已不再是边缘模型——因为它以 $V$ 为条件——但“边缘结构模型”这一名称仍被沿用。\mnote{参数 $\beta_3$ 一般没有“$V$ 的效应”这一因果解释。记住，我们假定的是治疗 $A$ 的可交换性、正值性与一致性，而不是性别 $V$ 的。}

我们可以用 IP 权重 $W^{A}$ 或 $SW^{A}$，通过加权最小二乘拟合线性回归模型 $\E[Y\given A,V]=\theta_0+\theta_1A+\theta_2VA+\theta_3V$ 来估计模型参数。在多数情形下，协变量向量 $L$ 应当包含 $V$。即便给定 $L$ 的其他成分时 $V$ 与 $A$ 独立、且并不需要 $V$ 来保证可交换性，把 $V$ 纳入 $L$ 一般也会提高边缘结构模型参数的估计效率。\mnote{代码：程序 12.6。}

由于我们考虑的是 $V$ 各水平内治疗效应的模型，现在可以选择在稳定化权重的分子中用 $f[A]$ 还是 $f[A\given V]$。基于稳定化权重 $SW^{A}(V)=f[A\given V]/f[A\given L]$ 的 IP 加权，一般会给出效应估计周围更窄的置信区间。对 $SW^{A}(V)$ 通常具有更高统计效率的一种直觉是：权重 $SW^{A}(V)$ 的方差小于权重 $SW^{A}$ 的方差。我们估计 $SW^{A}(V)$ 的方法与估计 $SW^{A}$ 相同，只是把 $V$ 加入权重分子所用的 logistic 模型。

研究者选择纳入边缘结构模型的 $L$ 的那个特定子集 $V$，应当只反映研究者的实质兴趣。例如，若研究者既相信 $V$ 可能是效应修饰因子，又更关心协变量 $V$ 各水平内的治疗因果效应而非整个总体中的效应，就应把 $V$ 纳入边缘结构模型。在我们的例子中，我们没有发现性别效应修饰的强证据，因为参数估计 $\hat\theta_2$ 周围的 95\% 置信区间为（$-2.2$, 1.9）。如果研究者选择把 $L$ 中所有变量都纳入边缘结构模型，则稳定化权重 $SW^{A}(L)$ 等于 1，IP 加权也就不必要了，因为在条件可交换性下，边缘结构模型此时就是那个用于完全校正 $L$ 所致混杂的（未加权）结局回归模型（见第 15 章）。出于这个原因，我们带着几分戏谑，把以可交换性所需的全部变量 $L$ 为条件的边缘结构模型称为\emph{虚假边缘结构模型}。\mnote{如果我们关心的是 2 个治疗 $A$ 与 $B$ 之间的交互作用（而不是变量 $V$ 对治疗 $A$ 的效应修饰；见第 5 章），我们会在边缘结构模型中同时纳入 $A$ 与 $B$ 的参数，并用两种治疗的联合概率作为分母来估计 IP 权重。我们会假定 $A$ 与 $B$ 的可交换性、正值性与一致性。}

在第一部分我们讨论过，效应修饰与混杂是两个逻辑上不同的概念。然而，许多学生难以理解这一区别，因为同样的统计方法——分层（第 4 章）或回归（第 15 章）——常常既用于校正混杂因子，又用于探查效应修饰。因此，用边缘结构模型讲授这些概念也许有些好处，因为此时校正混杂的方法（IP 加权）与探查效应修饰的方法（在边缘结构模型中加入治疗-协变量乘积项）是彼此分开的。

\section{删失与缺失数据}

在估计戒烟 $A$ 对体重增加 $Y$ 的因果效应时，我们把分析限制在 1982 年随访结束时测量了体重的 1566 个个体上。然而，另有 63 个个体符合我们的合格标准，却因 1982 年的体重未知而被排除在分析之外。只选择结局取值未缺失的个体——也就是把有缺失值的个体从分析中删失——可能引入选择偏倚，如第 8 章所讨论。

令删失 $C$ 表示 1982 年体重是否被测量：若体重未被测量（即个体被删失）取 1，若体重被测量（即个体未被删失）取 0。我们的分析必然被限制在未删失的个体（即 $C=0$ 者）上，因为只有他们的结局 $Y$ 取值已知。也就是说，在 12.2 节与 12.4 节，我们拟合的不是（加权）结局回归模型 $\E[Y\given A]=\theta_0+\theta_1A$，而是限制在 $C=0$ 个体上的模型 $\E[Y\given A,C=0]=\theta_0+\theta_1A$。

遗憾的是，即便在零假设下，当 $C$ 是治疗 $A$ 与结局 $Y$ 之间某条路径上的对撞节点、或是这样一个对撞节点的后代时，只选择未删失个体作分析也会引入偏倚。见图~8.3 至 8.6 的因果图。我们的数据与这些因果图所描绘的结构一致：治疗 $A$ 与删失 $C$ 相关——戒烟者中 5.8\% 被删失，未戒烟者中为 3.2\%；并且至少有一些 $Y$ 的预测因子与 $C$ 相关——被删失者的平均基线体重为 76.6 kg，未删失者为 70.8 kg。

由于失访所致的删失可能引入选择偏倚，我们一般关心的是“研究总体中无人被删失”时的因果效应。在我们的例子中，目标变为估计“人人都戒烟、且无人结局被删失”时的平均体重增加 $\E[Y^{a=1,c=0}]$，以及“人人都不戒烟、且无人结局被删失”时的平均体重增加 $\E[Y^{a=0,c=0}]$。于是，所关心的因果效应是 $\E[Y^{a=1,c=0}]-\E[Y^{a=0,c=0}]$，即 $A$ 与 $C$ 的联合效应，如第 8 章所讨论。使用上标 $c=0$ 明确表达了很多人提到治疗 $A$ 的因果效应时心中所想的那个因果对比，即便他们选择不用上标 $c=0$。

这一因果效应可以用 IP 权重 $W^{A,C}=W^{A}\times W^{C}$ 来估计，其中未删失个体的 $W^{C}=1/\Pr[C=0\given L,A]$，被删失个体的 $W^{C}=0$。在以下可识别条件下，IP 权重 $W^{A,C}$ 同时校正混杂与选择偏倚：给定 $L$ 时联合治疗 $(A,C)$ 的可交换性，即 $Y^{a,c=0}\indep(A,C)\given L$；$(A=a,C=0)$ 的联合正值性；以及一致性。如果 $L$ 中某些变量像图~8.4 那样受治疗 $A$ 影响，条件独立性 $Y^{a,c=0}\indep(A,C)\given L$ 一般不会成立。第三部分将说明，当 $L$ 的某些成分受治疗影响时，还有一些替代的可交换性条件，允许我们用 IP 加权估计 $A$ 与 $C$ 的联合效应。\mnote{删失与治疗的 IP 权重为 $W^{A,C}=1/f(A,C=0\given L)$，其中 $f(A,C=0\given L)=f(A\given L)\times\Pr[C=0\given L,A]$。}

记住，权重 $W^{C}=1/\Pr[C=0\given L,A]$ 构造的伪总体，其规模与删失之前的原研究总体相同，且其中从 $L$ 或 $A$ 指向 $C$ 的箭头都不存在。在我们的例子中，删失的 IP 权重 $W^{C}$ 的估计会构造出一个（近似）1566+63=1629 人的伪总体；在我们的假设下，其中没有选择偏倚，因为没有选择。也就是说，我们用权重 $W^{A,C}$ 拟合加权模型 $\E[Y\given A,C=0]=\theta_0+\theta_1A$，以估计整个总体中边缘结构模型 $\E[Y^{a,c=0}]=\beta_0+\beta_1a$ 的参数。\mnote{某些变量在 $f(A\given L)$ 的模型中系数可能为零，但在 $\Pr[C=0\given L,A]$ 的模型中不为零，反之亦然。大样本中，保留所有独立预测结局的变量 $L$ 总是更有效率。}

或者，也可以使用稳定化 IP 权重 $SW^{A,C}=SW^{A}\times SW^{C}$。删失权重 $SW^{C}=\Pr[C=0\given A]/\Pr[C=0\given L,A]$ 构造的伪总体，其规模与删失之后的原研究总体相同，且其中不存在从 $L$ 指向 $C$ 的箭头。在我们的例子中，删失的 IP 权重 $SW^{C}$ 的估计会构造出一个（近似）1566 名未删失个体的伪总体。也就是说，稳定化权重并不消除伪总体中的删失，而是使删失相对于已测协变量 $L$ 随机发生。因此，在“给定 $L$（与 $A$）时被删失者与未删失者可条件交换”这一假设下，伪总体中被删失个体的比例与研究总体中相同：有选择，但没有选择偏倚。\mnote{当 $\Pr[C=0\given A]$ 的模型正确时，$SW^{C}$ 的均值为 1。关于稳定化 IP 权重，见技术注记 12.2。代码：程序 12.7。估计的 $SW^{A,C}$ 从 0.35 到 4.09，均值为 1.00。}

为在我们的例子中得到 $\Pr[C=0\given L,A]$ 的参数估计，我们为未删失的概率拟合了一个 logistic 回归模型，数据为研究总体中的 1629 个个体。模型纳入的协变量与先前估计治疗权重时所用的相同。在这些参数限制下，我们对 1566 个未删失个体中的每一个都得到了 $\Pr[C=0\given L,A]$ 的估计 $\widehat{\Pr}[C=0\given L,A]$ 以及 $SW^{C}$ 的估计。使用稳定化权重 $SW^{A,C}=SW^{A}\times SW^{C}$，我们估计戒烟平均使体重增加 $\hat\theta_1=3.5$ kg（95\% 置信区间：2.5, 4.5）。这与我们先前用 IP 权重 $SW^{A}$ 得到的估计几乎相同，说明要么删失没有造成选择偏倚，要么我们已测的协变量无法消除这一偏倚。

\renewcommand{\TPname}{技术注记（Technical Point）}
\begin{techpoint}[关于稳定化权重的补充]
稳定化权重 $SW^{A}=f[A]/f[A\given L]$ 属于范围更大的一类稳定化权重 $g[A]/f[A\given L]$，其中 $g[A]$ 是 $A$ 的任一不是 $L$ 的函数。当使用非饱和结构模型时，权重 $g[A]/f[A\given L]$ 优于权重 $1/f[A\given L]$，因为存在一些函数 $g[A]$（通常 $f[A]$ 就是其中之一），可用来构造非饱和边缘结构模型中因果效应的更有效估计量。

尽管以权重 $g[A]/f[A\given L]$ 计算的 IP 加权均值 $\E\bigl[g(A)I(A=a)Y/f(A\given L)\bigr]$ 在可交换性与正值性下不再等于反事实均值 $\E[Y^{a}]$，但 IP 加权均值的 Hajek 形式
\[
\frac{\E\left[\dfrac{g(A)I(A=a)Y}{f(A\given L)}\right]}{\E\left[\dfrac{g(A)I(A=a)}{f(A\given L)}\right]}
\]
确实等于 $\E[Y^{a}]$，因为
\[
\E\Bigl[\frac{g(A)I(A=a)Y}{f(A\given L)}\Bigr]
= g(a)\,\E\Bigl[\frac{I(A=a)Y}{f(A\given L)}\Bigr]
= g(a)\,\E[Y^{a}],
\]
且 $\E\bigl[g(A)I(A=a)/f(A\given L)\bigr]=g(a)$。Hajek 均值是方程
\[
\E\left[\frac{g[A]}{f[A\given L]}(Y-u)\right]=0
\]
的解 $u$。类似地，在最简单的边缘结构模型 $\E[Y^{a}]=\beta_0+\beta_1a$ 中，以权重 $g[A]/f[A\given L]$ 计算的加权最小二乘估计量 $(\hat\beta_0,\hat\beta_1)$ 解如下估计方程
\[
\widehat{\E}\left[\frac{g[A]}{f[A\given L]}\bigl(Y-(\beta_0+\beta_1A)\bigr)\binom{1}{A}\right]=0.
\]
估计 $\hat\beta_0$ 是 $\E[Y^{0}]$ 的 Hajek 形式加权均值，$\hat\beta_0+\hat\beta_1$ 是 $\E[Y^{1}]$ 的 Hajek 形式加权均值，其中期望已换成样本平均。最后，仿照技术注记 2.2 的论证可以证明：在用权重 $g[A]/f[A\given L]$ 构造的伪总体中，给定 $A=a$ 时 $Y$ 的均值仍然等于 $\E[Y^{a}]$。
\end{techpoint}
\renewcommand{\TPname}{技术注记}

下面我们介绍 IP 加权之外另一种校正混杂与选择偏倚的方法：标准化。
